%% toc: P3 26 26.4
\chapter{Finite Signatures, Homogeneous Structures and the Types of the Random Graph}
\label{sec:p03-w03-p03-015-finite-signature}

The general finite-type criterion
(\cref{thm:p03-rn-characterization}) says that a complete theory in a
countable language with infinite models is \(\omega\)-categorical exactly
when it has only finitely many complete \(n\)-types for each \(n\).
That criterion places no restriction on the signature. This chapter
treats a narrower but frequently used route to its hypothesis: a
countably infinite structure in which every isomorphism between finitely
generated substructures extends to an automorphism. When the signature is
finite and has no function symbols, the finitely many atomic patterns of
an \(n\)-tuple already decide its complete type, so the type spaces are
finite and the criterion applies.

\emph{Inputs.} A finite signature with relation and constant symbols but
no function symbols, and a countably infinite ultrahomogeneous structure
in it; or, for the worked example, the graph axioms and extension axioms.
\emph{Conclusions.} Quantifier elimination, an explicit list of complete
\(n\)-types isolated by quantifier-free formulas, the exact number of
these types, and \(\omega\)-categoricity.

The hypothesis on function symbols is a correction. In his list of
corrigenda to \emph{Model Theory}, Hodges amends the finite-signature
application on p.~352 by adding the words ``with no function symbols''
after ``finite signature'' \cite[p.~3]{Hodges1993Corrigenda2011}. The same
list treats constant symbols and function symbols as separate kinds: its
correction to Corollary~2.7.6 on p.~72 states one clause for a signature
with a constant symbol and another for a finite signature with function
symbols \cite[p.~1]{Hodges1993Corrigenda2011}. Accordingly the theorem
below allows constant symbols. It is stated and proved here in its own
formulation; it is not a reading of the textbook page.
\Cref{thm:p03-rn-relational-application} and
\cref{thm:p03-fraisse54-categoricity-boundary} prove the purely relational
case without constants.

\section{Generated substructures and atomic patterns}
\label{sec:p03-w03-p03-015-setting}

\begin{hbdefinition}{def:p03-w03-p03-015-ultrahomogeneous}{Generated substructures, homogeneity}
Let \(L\) be a first-order signature and \(\mathcal M\) an
\(L\)-structure with domain \(M\). For \(X\subseteq M\), let
\(\langle X\rangle^{\mathcal M}\) be the least subset of \(M\) that
contains \(X\) and every constant \(c^{\mathcal M}\) and is closed under
every function \(F^{\mathcal M}\). It carries the restricted relations,
constants and functions. It is \emph{finitely generated} when \(X\) is
finite. When \(L\) has no constant symbols, \(\langle\varnothing\rangle^{\mathcal M}\)
is empty and the only map between two such sets is the empty map.

An \emph{isomorphism} \(f:\langle X\rangle^{\mathcal M}\to\langle Y\rangle^{\mathcal M}\)
is a bijection that sends each \(c^{\mathcal M}\) to \(c^{\mathcal M}\),
commutes with each \(F^{\mathcal M}\), and satisfies
\(R^{\mathcal M}(\bar u)\Leftrightarrow R^{\mathcal M}(f\bar u)\) for each
relation symbol \(R\) and tuple \(\bar u\) of the appropriate length.
The structure \(\mathcal M\) is \emph{ultrahomogeneous} when every
isomorphism between finitely generated substructures of \(\mathcal M\)
extends to an automorphism of \(\mathcal M\).
For a relational signature without constants this is the condition on
finite induced substructures used in \cref{thm:p03-fraisse54-graph}.
A \emph{Fra\"iss\'e limit} of a class of finitely generated
\(L\)-structures is a countable ultrahomogeneous structure whose finitely
generated substructures are, up to isomorphism, exactly the members of
the class.
\end{hbdefinition}

\begin{hbdefinition}{def:p03-w03-p03-015-patterns}{Atomic formulas and patterns on \(n\) slots}
Let \(L\) be finite, with relation symbols of positive finite arity and
constant symbols \(c_1,\ldots,c_k\) (\(k\geq0\)), and no function symbols.
Fix \(n\geq0\) and variables \(\bar x=(x_0,\ldots,x_{n-1})\).
The terms in these variables are just \(x_0,\ldots,x_{n-1},c_1,\ldots,c_k\).
Let \(\operatorname{At}_n\) be the set of atomic formulas
\(t=t'\) and \(R(t_1,\ldots,t_m)\) built from these terms; then
\[
 |\operatorname{At}_n|=(n+k)^2+\sum_{R}(n+k)^{\operatorname{ar}(R)}
\]
is finite. For an \(L\)-structure \(\mathcal M\) and
\(\bar a\in M^n\), the \emph{pattern} of \(\bar a\) is
\[
 \Delta^{\mathcal M}(\bar a)=\{\alpha\in\operatorname{At}_n:
              \mathcal M\models\alpha(\bar a)\}.
\]
For \(\Delta\subseteq\operatorname{At}_n\), its \emph{pattern formula} is
\[
 \pi_\Delta(\bar x)=\bigwedge_{\alpha\in\Delta}\alpha\ \land
   \bigwedge_{\alpha\in\operatorname{At}_n\setminus\Delta}\neg\alpha .
\]
Let \(P_n(\mathcal M)\) be the set of patterns of tuples in \(M^n\).
A tuple satisfies \(\pi_\Delta\) exactly when its pattern is \(\Delta\);
distinct pattern formulas are therefore jointly unsatisfiable, and
\(|P_n(\mathcal M)|\leq2^{|\operatorname{At}_n|}\).
For \(n=0\) the terms are the constants, the empty tuple has one
pattern, and an empty disjunction is falsity.
\end{hbdefinition}

\section{The finite-signature application}
\label{sec:p03-w03-p03-015-application}

\begin{hbtheorem}{thm:p03-w03-p03-015-finite-signature-application}{proved=annex:p03-w03-p03-015-app}{Finite signature without function symbols: homogeneity gives countable categoricity}
Let \(L\) be a finite signature with no function symbols; relation
symbols of positive finite arity and constant symbols are allowed. Let \(\mathcal M\) be a countably infinite
ultrahomogeneous \(L\)-structure and \(T=\operatorname{Th}(\mathcal M)\).
For every \(n\geq0\):
\begin{enumerate}
\item Two tuples in \(M^n\) with the same pattern lie in the same
      \(\operatorname{Aut}(\mathcal M)\)-orbit.
\item Every formula \(\varphi(\bar x)\) is equivalent modulo \(T\) to
      the disjunction of the \(\pi_\Delta\), \(\Delta\in P_n(\mathcal M)\),
      at whose realizations \(\varphi\) holds; so \(T\) has quantifier
      elimination.
\item Each \(\pi_\Delta\) with \(\Delta\in P_n(\mathcal M)\) isolates a
      complete type, and \(\Delta\mapsto\) that type is a bijection
      \(P_n(\mathcal M)\to S_n(T)\). Hence
      \(|S_n(T)|=|P_n(\mathcal M)|\leq2^{|\operatorname{At}_n|}\).
\item \(T\) is \(\omega\)-categorical: every countable model of \(T\) is
      isomorphic to \(\mathcal M\).
\end{enumerate}
In particular the conclusions hold for every infinite Fra\"iss\'e limit
in such a signature.
\end{hbtheorem}

The two hypotheses on \(L\) enter at two steps. Finitely many relation and
constant symbols, and no function symbols, make \(\operatorname{At}_n\)
finite. The absence of function symbols also makes
\(\langle\bar a\rangle^{\mathcal M}\) the finite set of values of the
terms \(x_i\) and \(c_j\), so a pattern describes the generated
substructure completely and an equal pattern defines an isomorphism.
Constant symbols cost nothing: they add \(k\) terms, and every isomorphism
of generated substructures fixes them. For a purely relational signature
without constants, existence of the limit of a Fra\"iss\'e class is
\cref{thm:p03-fraisse54-countable-language}.

\begin{hbtheorem}{thm:p03-w03-p03-015-function-boundary}{proved=annex:p03-w03-p03-015-fn}{The hypothesis on function symbols cannot be removed}
Let \(\mathcal Z=(\mathbb Z,s,t)\) with \(s(a)=a+1\) and \(t(a)=a-1\),
in the finite signature with two unary function symbols and nothing else.
Then \(\mathcal Z\) is countably infinite and ultrahomogeneous in the
sense of \cref{def:p03-w03-p03-015-ultrahomogeneous}, but
\(\operatorname{Th}(\mathcal Z)\) has infinitely many complete two-types
and is not \(\omega\)-categorical.
\end{hbtheorem}

Here the atomic formulas \(y=s^j(x)\), \(j\geq0\), are infinitely many,
which is exactly the step that fails. Finiteness of the signature is also
needed: with infinitely many unary predicates there is an ultrahomogeneous
countable structure whose theory is not \(\aleph_0\)-categorical
(\cref{thm:p03-fraisse54-categoricity-boundary}).

\section{The random graph: the complete type calculation}
\label{sec:p03-w03-p03-015-random-graph}

Let \(L_E=\{E\}\) with \(E\) binary, and let \(T_{\rm RG}\) be the theory
of \cref{def:p03-fraisse54-graph-axioms}: irreflexivity, symmetry, and for
all \(r,s\geq0\) the extension axiom asserting that for any \(r+s\)
distinct vertices there is a further vertex adjacent to the first \(r\)
and to none of the last \(s\). Its countable model \(\mathcal R\) is the
limit of the class of finite simple graphs (\cref{thm:p03-fraisse54-graph},
\cref{thm:p03-random-graph}). No probability distribution is involved.

\begin{hbdefinition}{def:p03-w03-p03-015-graph-patterns}{Graph patterns}
A \emph{graph pattern} on \(n\) slots is a pair \((\varepsilon,\eta)\)
where \(\varepsilon\) is an equivalence relation on
\(\{0,\ldots,n-1\}\) and \(\eta\) is an irreflexive symmetric relation on
the set of its classes, that is, a simple graph whose vertices are the
\(\varepsilon\)-classes. Write \([i]\) for the class of \(i\). Its formula is
\[
 \begin{aligned}
 \gamma_{\varepsilon,\eta}(\bar x)={}&
 \bigwedge_{i\,\varepsilon\,j}x_i=x_j\ \land
 \bigwedge_{\neg\,i\,\varepsilon\,j}x_i\neq x_j\\
 &\land\bigwedge_{[i]\,\eta\,[j]}E(x_i,x_j)\ \land
 \bigwedge_{[i]\neq[j],\ \neg[i]\,\eta\,[j]}\neg E(x_i,x_j),
 \end{aligned}
\]
with \(i,j\) ranging over \(\{0,\ldots,n-1\}\). Let \(\Gamma_n\) be the
set of graph patterns on \(n\) slots. For \(n=0\) there is one pattern,
with formula \(\top\). Every tuple of vertices of a simple graph
satisfies exactly one \(\gamma_{\varepsilon,\eta}\): its equality
partition together with the adjacency of its distinct entries.
\end{hbdefinition}

\begin{hblemma}{thm:p03-w03-p03-015-graph-realization}{proved=annex:p03-w03-p03-015-gr}{Realizing graph patterns}
Every model of \(T_{\rm RG}\) realizes every graph pattern on \(n\)
slots, for every \(n\geq0\).
\end{hblemma}

\begin{hblemma}{thm:p03-w03-p03-015-graph-partial}{proved=annex:p03-w03-p03-015-gp}{Partial maps preserve formulas}
Let \(\mathcal G,\mathcal H\models T_{\rm RG}\). Call an injective map
\(f\) from a finite subset of \(G\) into \(H\) a \emph{finite partial
isomorphism} when \(uE^{\mathcal G}v\Leftrightarrow f(u)E^{\mathcal H}f(v)\)
for all \(u,v\) in its domain. For every such \(f\), every formula
\(\varphi(\bar x)\) and every tuple \(\bar a\) from the domain of \(f\),
\[
 \mathcal G\models\varphi(\bar a)\iff\mathcal H\models\varphi(f\bar a).
\]
\end{hblemma}

\begin{hbtheorem}{thm:p03-w03-p03-015-graph-types}{proved=annex:p03-w03-p03-015-gt}{The complete \(n\)-types of the random graph}
For every \(n\geq0\):
\begin{enumerate}
\item tuples in any two models of \(T_{\rm RG}\) with the same graph
      pattern satisfy the same formulas; in particular \(T_{\rm RG}\) is
      complete;
\item every formula \(\varphi(\bar x)\) is equivalent modulo
      \(T_{\rm RG}\) to a disjunction of formulas
      \(\gamma_{\varepsilon,\eta}\), so \(T_{\rm RG}\) has quantifier
      elimination;
\item each \(\gamma_{\varepsilon,\eta}\) isolates a complete type, and
      \((\varepsilon,\eta)\mapsto\) that type is a bijection
      \(\Gamma_n\to S_n(T_{\rm RG})\);
\item the number of complete \(n\)-types is
      \[
       |S_n(T_{\rm RG})|=|\Gamma_n|
        =\sum_{k=0}^{n}{n\brace k}\,2^{k(k-1)/2},
      \]
      where \({n\brace k}\) counts partitions of an \(n\)-set into \(k\)
      nonempty blocks and \({0\brace0}=1\); for \(n=0,1,2,3,4\) this is
      \(1,1,3,15,127\);
\item the \(\operatorname{Aut}(\mathcal R)\)-orbits on \(n\)-tuples of
      \(\mathcal R\) are exactly the sets of realizations of the
      \(\gamma_{\varepsilon,\eta}\); there are \(|\Gamma_n|\) of them.
\end{enumerate}
\end{hbtheorem}

\begin{hbexample}{def:p03-w03-p03-015-small-types}{Two and three variables}
For \(n=2\) the three types are isolated by \(x_0=x_1\),
\(x_0\neq x_1\land E(x_0,x_1)\) and \(x_0\neq x_1\land\neg E(x_0,x_1)\).
For \(n=3\): one type has all three entries equal; for each of the three
ways to make exactly two entries equal, the two distinct values are
adjacent or not, giving six types; three distinct entries carry any of the
\(2^3=8\) labelled graphs on three vertices. The total is \(1+6+8=15\).
The general bound \(2^{|\operatorname{At}_n|}=2^{2n^2}\) of
\cref{def:p03-w03-p03-015-patterns} is far larger; the graph axioms and
the laws of equality leave only the graph patterns.
\end{hbexample}

\begin{hbtheorem}{thm:p03-w03-p03-015-random-graph-instance}{proved=annex:p03-w03-p03-015-rg}{The random graph as an instance of the finite-signature application}
In \cref{thm:p03-w03-p03-015-finite-signature-application} take
\(L=L_E\), which is finite with no function or constant symbols, and
\(\mathcal M=\mathcal R\). The map sending a pattern
\(\Delta\in P_n(\mathcal R)\) to its equality partition and the adjacency
of its classes is a bijection \(P_n(\mathcal R)\to\Gamma_n\), and
\(\pi_\Delta\) is equivalent modulo \(T_{\rm RG}\) to the corresponding
\(\gamma_{\varepsilon,\eta}\). The theorem therefore returns the same
types and the same count \(|S_n(\operatorname{Th}(\mathcal R))|=|\Gamma_n|\),
and \(\operatorname{Th}(\mathcal R)\) is the set of consequences of
\(T_{\rm RG}\) and is \(\omega\)-categorical.
\end{hbtheorem}

The two computations use different inputs. \Cref{thm:p03-w03-p03-015-graph-types}
uses only the axioms: the extension axioms realize every pattern and
extend every finite partial isomorphism by one point, inside any model.
\Cref{thm:p03-w03-p03-015-random-graph-instance} uses ultrahomogeneity of
the one countable model; the extension axioms then enter only through the
realization of every graph pattern in \(\mathcal R\).

\begin{hbtheorem}{thm:p03-w03-p03-015-named-vertex}{proved=annex:p03-w03-p03-015-nv}{A constant symbol: the random graph with a named vertex}
Let \(v\) be a vertex of \(\mathcal R\) and let \((\mathcal R,v)\) be the
\(\{E,c\}\)-structure with \(c\) interpreted by \(v\). Then
\((\mathcal R,v)\) is ultrahomogeneous, so
\cref{thm:p03-w03-p03-015-finite-signature-application} applies with one
constant symbol. Its complete \(n\)-types correspond to the graph patterns
of the \((n+1)\)-tuples \((\bar a,v)\), and
\[
 |S_n(\operatorname{Th}(\mathcal R,v))|=|\Gamma_{n+1}|,
\]
so for \(n=0,1,2,3\) the counts are \(1,3,15,127\).
In particular \(\operatorname{Th}(\mathcal R,v)\) is \(\omega\)-categorical.
\end{hbtheorem}

\section{An effective quantifier-elimination procedure}
\label{sec:p03-w03-p03-015-effective-elimination}

The type calculation also gives a procedure for answering a formula query
from the equality and adjacency data of its free variables. The procedure
below constructs the disjunction in
\cref{thm:p03-w03-p03-015-graph-types}; it does not presuppose an oracle
for truth in the random graph. Its input is a finite first-order formula
in the language \(\{E,=\}\), with no named parameters. The extension
axioms, rather than a sampled finite graph, justify its quantifier steps
\cite{Fraisse1954,Rado1964}.

\begin{hbmethod}{def:p03-w03-p03-015-elimination}{Compute the satisfying graph patterns}
Work with \(\neg,\land,\exists\); express the other Boolean connectives
and \(\forall\) using these. Rename bound variables to avoid capture.
For a list \(\bar x=(x_0,\ldots,x_{n-1})\) of distinct variables
containing every free variable of \(\varphi\), compute a subset
\(Q_{\bar x}(\varphi)\) of \(\Gamma_n\). The list may contain variables
that do not occur freely in the formula. In a recursive call below a
quantifier, choose its bound variable \(y\) fresh for this list.

For a pattern \(\rho'\in\Gamma_{n+1}\), let
\(\rho'\mathbin{\upharpoonright}n\) be its restriction to the first
\(n\) slots: restrict the equality partition and the graph to the
classes meeting those slots. The following clauses are instructions on
finite sets:
\begin{enumerate}
\item For \(x_i=x_j\), retain exactly the patterns whose partition puts
      \(i,j\) in the same class. For \(E(x_i,x_j)\), retain exactly the
      patterns in which their classes are distinct and adjacent.
      Truth and falsity retain \(\Gamma_n\) and the empty set,
      respectively.
\item Complement the computed set for negation, relative to
      \(\Gamma_n\); intersect the two computed sets for conjunction.
      Both operands use the same variable list, including its unused slots.
\item For \(\exists y\,\psi(\bar x,y)\), first compute
      \(Q_{\bar x,y}(\psi)\). Return the set of restrictions
      \(\rho'\mathbin{\upharpoonright}n\) of its members, removing
      repetitions. Thus
      \[
       Q_{\bar x}(\exists y\,\psi)
       =\{\rho'\mathbin{\upharpoonright}n:
                    \rho'\in Q_{\bar x,y}(\psi)\}.
      \]
\item Return the quantifier-free formula
      \[
       \bigvee_{\rho\in Q_{\bar x}(\varphi)}\gamma_\rho(\bar x),
      \]
      where \(\gamma_\rho\) is the graph-pattern formula of
      \cref{def:p03-w03-p03-015-graph-patterns}. An empty disjunction
      denotes \(\bot\), and an empty conjunction denotes \(\top\).
\end{enumerate}
To enumerate \(\Gamma_n\), enumerate the partitions of the finite slot
set and then every simple graph on each partition's classes. Label the
classes by their least slots, so comparison and duplicate removal are
finite syntactic operations.
\end{hbmethod}

For an existential step one can instead extend each candidate pattern
directly. If it has \(k\) equality classes, the new slot can join one of
the \(k\) existing classes, with adjacency then forced, or form a new
class with any of the \(2^k\) adjacency choices to the old classes.
These \(k+2^k\) possibilities exhaust its one-slot extensions. The
first group represents old witnesses and the second represents fresh
ones. Both groups are needed even though the extension axioms themselves
provide fresh vertices.

\begin{hbtheorem}{thm:p03-w03-p03-015-effective-elimination}{proved=annex:p03-w03-p03-015-effective-elimination}{Termination, elimination and decision in the random-graph theory}
For every finite \(\{E,=\}\)-formula \(\varphi(\bar x)\), the procedure
of \cref{def:p03-w03-p03-015-elimination} terminates. For every
\(\mathcal G\models T_{\rm RG}\) and tuple \(\bar a\in G^n\),
\[
 \mathcal G\models\varphi(\bar a)
 \quad\Longleftrightarrow\quad
 \text{the graph pattern of \(\bar a\) belongs to
       \(Q_{\bar x}(\varphi)\)}.
\]
Consequently its output is equivalent to \(\varphi\) in every model of
\(T_{\rm RG}\). A formula is satisfiable with \(T_{\rm RG}\) exactly
when its computed set is nonempty, and its universal closure is a
consequence of \(T_{\rm RG}\) exactly when that set is \(\Gamma_n\).
For a sentence, \(\Gamma_0\) has one member, so this is a decision
procedure for the complete theory of the random graph. It makes no
polynomial-time or polynomial-output-size claim.
\end{hbtheorem}

\begin{hbexample}{def:p03-w03-p03-015-elimination-examples}{A distinguishing neighbour and an old witness}
Consider
\(\exists y(E(x_0,y)\land\neg E(x_1,y))\).
If the two free slots coincide, either conjunct contradicts the other,
so no extension pattern survives. If they are distinct, a fresh class
adjacent to the class of \(x_0\) and not to that of \(x_1\) survives,
regardless of the edge between those two classes. The output is therefore
equivalent to \(x_0\neq x_1\). The extension axiom realizes that fresh
pattern over every pair of distinct vertices, not merely over some pair.

For \(\exists y(y=x_0)\), the extension that joins the old class of
\(x_0\) survives, so the output is \(\top\). A procedure considering
fresh witnesses only would incorrectly return \(\bot\). Conversely,
\(\exists y(y\neq x_0\land E(x_0,y))\) requires a fresh adjacent
class and is also equivalent to \(\top\).

Finally, \(\forall y\,E(x_0,y)\) is equivalent to \(\bot\): the
old-class extension \(y=x_0\) witnesses its negation, since the graphs
are irreflexive. These examples use both equality and nonadjacency;
recording only edges between distinct variable names would lose that
information when two names denote the same vertex.
\end{hbexample}

The finite enumeration grows with the active variable list: its size is
the sum of partition and graph counts in
\cref{thm:p03-w03-p03-015-graph-types}. A quantifier step can compute
its image by examining each retained pattern on one more slot once.
This is a concrete termination argument, not a bound on the complete
formula-processing cost. The method decides truth relative to
\(T_{\rm RG}\); finite graphs generally fail its extension axioms and
are not the models to which the theorem applies.


\section{Limitations and relationships}
\label{sec:p03-w03-p03-015-limits}

The type counts \(|\Gamma_n|\) grow without bound in \(n\); the criterion
needs finiteness for each \(n\), not a bound uniform in \(n\). For a
finite structure the type spaces are also finite but the theory has no
infinite model (\cref{thm:p03-rn-finite-boundary}), which is why the
application assumes \(\mathcal M\) infinite. A function symbol can make
the atomic formulas in two variables infinite even in a finite signature
(\cref{thm:p03-w03-p03-015-function-boundary}), and infinitely many
relation symbols can defeat categoricity even for an ultrahomogeneous
limit (\cref{thm:p03-fraisse54-categoricity-boundary}). The dense linear
order gives a parallel calculation with ordered partitions in place of
graph patterns (\cref{thm:p03-rn-dlo-type-count}). The equality of
\(n\)-types and automorphism orbits in the countable model also follows
from \cref{thm:p03-rn-orbits}.
